\begin{table}[H]
    \centering
    \renewcommand{\arraystretch}{1.3}
    \caption{Core Functions of the TrackCircuitBranch Module}
    \begin{tabular}{p{5.2cm} p{10.2cm}}
        \toprule
        \textbf{Function} & \textbf{Description} \\
        \midrule
        enforceTrackSegment \newline OccupancyInterlocking & Main method triggered upon occupancy transition. Identifies protecting signals and initiates enforcement. \\
        checkTrackSegmentExists & Verifies the existence of the track segment using the database. \\
        checkTrackSegmentActive & Checks if the track segment is operational and enabled. \\
        getTrackSegmentState & Fetches current metadata: occupancy, activity, assigned status, and protecting signals. \\
        getProtectingSignals \newline FromBothSources & Combines protecting signal IDs from both DB relations and configuration metadata. \\
        getProtectingSignals \newline FromDatabase & Queries signal-track protection mapping from the relational table. \\
        getProtectingSignals \newline FromTrackSegmentData & Extracts protecting signal array from track segment's own data field. \\
        enforceSignalToRed & Attempts to set a single signal to RED; verifies status and logs failures. \\
        enforceMultipleSignalsToRed & Batch-enforces RED on all protecting signals and reports any failures. \\
        verifySignalIsRed & Confirms a signal's current aspect is RED. \\
        areAllSignalsAtRed & Checks if all listed signals are already at RED state. \\
        handleInterlockingFailure & Logs failure, emits freeze signal, and triggers interlockingFailure signal. \\
        emitSystemFreeze & Issues a freeze event with detailed error reasons and affected entities. \\
        logCriticalFailure & Logs a detailed system-level failure with timestamp and thread. \\
        formatFailureDetails & Assembles a detailed failure message including time, signal IDs, and errors. \\
        \bottomrule
    \end{tabular}
    \label{tab:track-circuit-branch-functions}
\end{table}
